\section{Introduction}

Across cryptocurrency ecosystems, illicit actors exploit pseudonymous identities, rapid cross-border value transfer, and composable blockchain services to conduct scams, thefts, money laundering, ransomware payments, sanctions-related value transfer, and other high-risk financial activities. These activities threaten user assets, exchange compliance, and the credibility of decentralized finance infrastructures. The scale of the problem is already substantial. TRM Labs reported that adjusted incoming value to illicit cryptocurrency wallets reached USD 158 billion in 2025, the highest level observed in the previous five years \cite{trmlabs2026crypto}. The FBI Internet Crime Complaint Center reported that U.S. complaints involving cryptocurrency caused more than USD 11 billion in losses in 2025, with 181,565 cryptocurrency-related complaints \cite{fbi2026crypto}. Chainalysis further reported that known illicit cryptocurrency addresses received USD 40.9 billion in 2024, and this number is expected to increase as attribution improves \cite{chainalysis2025crypto}. These observations show that blockchain fraud detection is not only a technical classification task, but also a large-scale risk-control problem with direct financial and societal impact.

The difficulty of blockchain fraud detection comes from the way fraudulent behavior is organized on-chain. Risk is rarely determined by a single transaction amount, a single address feature, or a local attribute in isolation. Fraudulent strategies often unfold through fund flows, counterparties, repeated interactions, transaction motifs, multi-hop propagation paths, and coordinated activities among wallets, exchanges, contracts, mixers, and services. This makes graph learning a natural modeling choice. At the transaction level, nodes can represent transactions and edges can represent transaction-flow relations, allowing graph models to capture local transfer structures and suspicious flow patterns. At the actor or address level, nodes can represent blockchain entities and edges can describe transfer or interaction relations, allowing graph models to capture longer-term behavioral dependencies. Compared with feature-only classifiers, graph learning provides a way to use relational evidence that is difficult to observe from tabular features alone, especially when fraud signals are sparse, distributed, and intentionally obfuscated.

The temporal nature of blockchain systems further complicates this graph-learning problem. Transactions arrive continuously, addresses and actors change their activity patterns, and fraud strategies evolve in response to regulation, exchange compliance, market cycles, and deployed detection systems. A pattern that indicates suspicious activity in an early time window may weaken, move to another service, or reappear as a modified strategy in later windows. Prior real-time fraud detection research has emphasized that temporal causality must be respected, because using future graph information to predict current risk creates temporal leakage \cite{lu2022bright}. Dynamic graph learning therefore appears to be a relevant direction for blockchain fraud detection. Existing dynamic graph methods commonly use temporal neighborhoods, historical event sequences, recurrent updates, memory modules, or cumulative graph states to model evolving structures and node representations \cite{zheng2024dynamicgnnsurvey}. However, these mechanisms often assume that historical events, historical neighborhoods, or historical representations can be retained and revisited during later training stages.

Such assumptions are difficult to satisfy in practical blockchain risk-control systems. A full historical blockchain graph grows continuously as new transactions, addresses, edges, and multi-hop neighborhoods appear. Keeping this cumulative structure for model updates can introduce substantial storage, GPU memory, and computation costs. Reconstructing historical neighborhoods during online updates can also introduce latency, which is undesirable for fraud screening systems that require timely risk scoring. In addition, temporal evaluation must avoid future-information leakage: when the model predicts risk in a current time window, it should not use graph structures that only become visible in later windows. These constraints motivate the \emph{task-only snapshot} setting studied in this paper. The temporal blockchain fraud graph is divided into a chronological sequence of graph snapshots. At each task, the model can use only the snapshot constructed within the current time window. After moving to later tasks, it cannot store, access, or replay historical nodes, historical edges, adjacency matrices, or historical subgraphs.

Figure~\ref{fig:full-history-graph} illustrates the deployment burden of maintaining a full historical graph for online blockchain fraud detection.

\begin{figure}[t]
    \centering
    \includegraphics[width=\linewidth]{figures/full_history_graph.png}
    \caption{Industrial motivation of full-history temporal graph maintenance. Continuously maintaining a cumulative blockchain graph requires storing historical nodes, edges, adjacency matrices, and multi-hop neighborhoods, which increases storage, memory, computation, and latency in online fraud detection.}
    \label{fig:full-history-graph}
\end{figure}

The central problem under this setting is how to adapt to newly emerging blockchain fraud patterns while maintaining fraud-discriminative semantics learned from previous temporal tasks. Training only on the current snapshot may overfit short-term local behaviors and weaken stable risk knowledge learned from earlier windows. Storing the complete historical graph would provide replay information, but this violates the task-only constraint. We therefore focus on preserving structure-free semantic information that remains useful across time. In blockchain fraud detection, this information can be viewed as two complementary parts: normal behavioral semantics that describe legitimate transaction and interaction patterns, and abnormal fraud semantics that describe high-risk fund flows, suspicious counterparties, and fraud-related structural evidence.

To address this problem, we propose \textbf{TASD-CL}, a Temporal Adaptive Semantic Decomposition framework for task-only blockchain graph fraud detection. TASD-CL explicitly decomposes node representations into a normal semantic subspace and an abnormal semantic subspace, and then adaptively combines them through alpha-guided routing for risk-aware message aggregation and fraud prediction. This semantic decomposition gives the model a clear interface for continual preservation. Instead of storing historical graph structures, TASD-CL retains lightweight structure-free information that summarizes or protects fraud-discriminative semantics across temporal tasks.

Implementing this idea raises several challenges. The first challenge is structure unavailability. Once the model moves to a new task, historical nodes, edges, adjacency matrices, and neighborhood subgraphs cannot be used as replay data. TASD-CL addresses this challenge by preserving only structure-free information, including semantic parameter importance, compact subspace prototypes, and previous model parameters. The second challenge is semantic drift. Normal blockchain behavior may change with market cycles and application usage, while fraud strategies may migrate across services or adopt new laundering patterns. Ordinary graph aggregation can entangle normal interaction semantics and abnormal risk semantics, making them difficult to preserve separately. TASD-CL addresses this challenge by separating normal and abnormal semantic subspaces and using alpha-guided routing to keep their risk roles explicit during message passing. The third challenge is severe class imbalance. Fraudulent nodes are much rarer than normal nodes, so abnormal semantics can be overwhelmed by majority normal patterns during sequential training. TASD-CL mitigates this challenge through task-adaptive focal supervision, abnormal-aware semantic distillation, and prototype extraction safeguards that avoid unreliable minority prototypes. The fourth challenge is cross-granularity modeling. Blockchain fraud graphs may be built at transaction level or actor/address level, and their node and edge meanings differ. TASD-CL handles both cases as temporal node risk recognition over evolving blockchain interaction graphs, while its semantic decomposition operates on learned node representations rather than dataset-specific node definitions.

Built on this design, TASD-CL introduces three lightweight semantic preservation components. Subspace-Stratified Fisher regularization (SSF) stabilizes semantically important parameters according to their roles in decomposition and routing. Subspace Prototype Condensation (SPC) summarizes historical normal and abnormal semantic distributions through compact Gaussian prototypes without storing historical graph structures. Subspace-Conditioned Distillation (SCD) selectively distills high-confidence normal and abnormal semantics from the previous model on the current snapshot. Together, these components allow TASD-CL to adapt to new temporal snapshots while preserving fraud-discriminative semantics under task-only constraints.

The main contributions of this paper are summarized as follows.

\begin{itemize}
    \item We formulate \emph{task-only snapshot continual blockchain graph fraud detection}, where each temporal task can use only the current blockchain graph snapshot and future tasks cannot access or replay historical graph structures.

    \item We design a temporal adaptive semantic decomposition framework that explicitly separates node representations into normal and abnormal semantic subspaces and performs risk-aware prediction through alpha-guided routing.

    \item We develop a three-level semantic preservation architecture for TASD-CL. SSF preserves fraud-discriminative semantics at the parameter level, SPC preserves them at the representation level through structure-free semantic prototypes, and SCD preserves them at the semantic level through subspace-conditioned distillation.

    \item We evaluate TASD-CL on two complementary temporal blockchain fraud datasets, Elliptic for transaction-level risk recognition and Elliptic++ Actor for actor/address-level risk recognition. Experiments using AUC-ROC and MacroF1 show the effectiveness of TASD-CL under severe class imbalance and task-only temporal snapshots.
\end{itemize}